\documentclass[11pt]{article}
\usepackage{icmlko}
\begin{document}

\title{육십다섯 노트 동안 죽어 있던 손잡이 --- 깊이를 올린 날 되살아났는데 아무도 안 돌아봤다}
\author{월드모델 실험실}
\date{노트 519 · 2026년 8월 4일}
\maketitle

\begin{abstract}
노트 516 이 \texttt{TRAINR}(나이 지수 감쇠 무게)을 채택했다. 원장을 뒤지다
같은 손잡이가 \textbf{이미 반증돼 있는 것}을 찾았다.

\begin{quote}
\textbf{노트 451 마무리 --- ``무게로는 못 산다. 나이 지수 감쇠 무게는 어떤
$\tau$ 도 균등을 못 이긴다: $\tau{=}8$ $-0.0005$ · $\tau{=}4$ $-0.0038$ ·
$\tau{=}2$ $-0.0130$ · $\tau{=}1$ $-0.0294$. 단조로 나빠진다.''}
\end{quote}

\textbf{같은 손잡이가 지금 판 $+0.0036$ · KR 만화 $+0.031$ 을 산다.}
무엇이 달라졌는지 후보를 하나씩 떨어뜨렸다.

\begin{center}
\begin{tabular}{ll}
\hline
후보 & 결과 \\
\hline
\texttt{TRAINCAP} $\cdot$ \texttt{TRAINW} 가 되살렸나 & \textbf{✗} 세 단계 다 양수($+0.0023 / +0.0028 / +0.0028$) \\
격자가 봉우리 왼쪽에만 있었나 & \textbf{✗} 반감기 $5.55$ 도 $+0.0022$ \\
\texttt{DOMDROP} 이 되살렸나 & \textbf{✗} 끄면 오히려 $+0.0044$ \\
무게 재정규화 차이인가 & \textbf{✗} 차 $-0.0001$ \\
\hline
\textbf{깊이} & \textbf{✓} \\
\hline
\end{tabular}
\end{center}

\begin{center}
\begin{tabular}{lrrrl}
\hline
& 균등 & $\tau{=}8$ & 차 & 부호 \\
\hline
깊이 $6$(노트 451 당시) & $+0.4839$ & $+0.4840$ & $+0.0001$ & $6/12$ \\
깊이 $12$(지금) & $+0.4812$ & $+0.4848$ & $\mathbf{+0.0036}$ & $\mathbf{12/12}$ \\
\hline
\end{tabular}
\end{center}

\textbf{노트 451 은 그때 맞았다.} 깊이 $6$ 에서 그 손잡이는 정말 죽어
있었고($+0.0001$, 그때 잰 $-0.0005$ 와 잡음 안에서 같다). 노트 457 이
깊이를 $12$ 로 올렸고, \textbf{그날 이 갈래가 살아났는데 아무도 안
돌아봤다.}
\end{abstract}

\begin{figure}[t]
\centering
\includegraphics[width=\textwidth]{figs/f1.pdf}
\caption{\textbf{가} --- 같은 손잡이 다른 용량. \textbf{나} --- 후보 다섯.
\textbf{다} --- 죽고 되살아나기까지.}
\end{figure}

\section{조항을 넓힌다}

노트 456 이 적었다 --- \textbf{``자가 바뀌면 그 자로 닫았던 갈래는 다시
연다.''} 그때는 \emph{채점하는 자}(집 밖 짝이 하나에서 둘로) 이야기였다.

\begin{quote}
\textbf{조항(넓힘)} --- \textbf{용량이 바뀌면 그 용량에서 닫았던 갈래도
다시 연다.} 깊이 $\cdot$ \texttt{it} $\cdot$ $K$ 를 바꾸는 채택은
\emph{그 이전에 죽은 자료 쪽 후보 목록}을 함께 남긴다.
\end{quote}

노트 406 이 이미 같은 모양을 봤다 --- ``죽은 손잡이는 자료가 자라면
되살아난다''(노트 272 가 죽인 깊이가 자료가 늘자 $4 \to 6$ 으로 살아났다).
거기서는 \textbf{자료}가 \textbf{용량}을 되살렸고, 여기서는
\textbf{용량}이 \textbf{자료 쓰는 법}을 되살렸다. \textbf{손잡이들은 서로의
사망 원인이자 부활 원인이다.}

\section{얼마를 흘렸나}

\begin{center}
\begin{tabular}{lr}
\hline
& 판 \\
\hline
노트 $457$ 이 깊이 $12$ 로 산 것 & $-0.0032$(집 밖으로 삼) \\
그날 같이 살아난 \texttt{TRAINR} & $\mathbf{+0.0036}$ \\
\hline
\end{tabular}
\end{center}

\textbf{깊이 $12$ 는 판을 잃는 채택이었는데, 그것이 연 문 뒤에 판
$+0.0036$ 이 있었다.} 노트 $457$ 부터 노트 $516$ 까지 육십다섯 노트 동안
그 문이 닫혀 있었다.

\section*{바로 해야 할 일}

이 사이클에서 용량은 안 바꿨지만 \textbf{자료 쓰는 법을 네 번 바꿨다}
(\texttt{DOMDROP} $\cdot$ \texttt{TRAINCAP} $\cdot$ \texttt{TRAINW}
$\cdot$ \texttt{TRAINR}). 조항을 대칭으로 읽으면 ---

\begin{quote}
\textbf{자료 쓰는 법이 바뀌었으니 그 이전에 죽은 \emph{용량} 후보를 다시
연다.}
\end{quote}

죽은 용량 후보: 깊이 $16 \cdot 24$(노트 458), \texttt{it} $120 \cdot 400$
(노트 $459 \sim 461 \cdot 473$), $K{=}64$(노트 456), $lr$ $0.06 \cdot
0.16$(노트 459 $\cdot$ 474). \textbf{다음 사이클의 첫 일이다.}

\section*{남은 것}

\begin{itemize}
\item \textbf{죽은 용량 후보 여섯을 지금 챔피언에서 다시} --- 위.
\item \textbf{다른 죽은 자료 후보} --- 원장에 ``반증''으로 닫힌 것이
스물 남짓이다. 그중 용량 감응인 것이 더 있나.
\item \textbf{조항을 기계로} --- 채택할 때마다 ``이 채택이 되살릴 수 있는
갈래''를 원장에 자동으로 적게 하면 사람이 안 잊는다.
\end{itemize}

\end{document}
